% !TEX root = master.tex
% !TEX program = lualatex

\chapter{Pointers - 4}
\label{chap:pointers_}

\section{Pointeurs et tableaux}

Nous allons clarifier le lien fondamental entre pointeurs et tableaux.

On peut écrire :

\begin{lstlisting}[language=C]
double* ptr = calloc(3, sizeof(double));
\end{lstlisting}

Ici :

\begin{itemize}
\item ptr est un pointeur vers double.
\item On alloue dynamiquement de la place pour 3 doubles.
\item La mémoire est initialisée à 0.
\end{itemize}

Important :

\[
\boxed{
*ptr \text{ lit uniquement le premier élément}
}
\]

Même si la zone allouée contient 3 doubles,
\texttt{*ptr} accède seulement au premier.

\section{Accès via la syntaxe tableau}

On peut accéder aux éléments ainsi :

\begin{lstlisting}[language=C]
ptr[0]
ptr[1]
ptr[2]
\end{lstlisting}

Rappel fondamental :

\[
\boxed{
*ptr \equiv ptr[0]
}
\]

Plus généralement :

\[
ptr[i] \equiv *(ptr + i)
\]

La notation tableau est donc une simple syntaxe
pour l’arithmétique des pointeurs.

\section{Différence fondamentale : pointeur vs tableau}

Considérons maintenant :

\begin{lstlisting}[language=C]
double tab[3];
\end{lstlisting}

Différence essentielle :

\begin{itemize}
\item Un pointeur est une variable contenant une adresse.
\item Un tableau est une zone mémoire continue.
\item Le nom du tableau n’est pas une variable contenant une adresse.
\end{itemize}

Un tableau :

\begin{itemize}
\item est alloué statiquement (souvent sur la stack),
\item ne possède pas de case mémoire supplémentaire pour stocker une adresse,
\item ne peut pas être réaffecté.
\end{itemize}

On peut voir un tableau comme :

\[
\boxed{
T[N] \approx T* const
}
\]

C’est un pointeur constant vers une zone mémoire continue.

\section{Mémoire : stack vs heap}

Rappel organisation mémoire :

\begin{itemize}
\item Stack : variables locales, tableaux statiques.
\item Heap : mémoire dynamique via malloc/calloc.
\end{itemize}

Exemple :

\begin{lstlisting}[language=C]
double P1[N][M];
double* P2[N];
double** P3;
\end{lstlisting}

\section{Cas 1 : tableau bidimensionnel}

\begin{lstlisting}[language=C]
double P1[N][M];
\end{lstlisting}

Caractéristiques :

\begin{itemize}
\item Alloué entièrement sur la stack.
\item Mémoire continue.
\item Contient N × M doubles consécutifs.
\end{itemize}

Accès :

\begin{lstlisting}[language=C]
P1[i][j]
\end{lstlisting}

Mémoire contiguë.

\section{Cas 2 : tableau de pointeurs}

\begin{lstlisting}[language=C]
double* P2[N];
\end{lstlisting}

Ici :

\begin{itemize}
\item P2 est un tableau de N pointeurs.
\item Chaque P2[i] peut pointer vers une zone différente.
\item Les lignes peuvent avoir des tailles différentes.
\end{itemize}

Allocation typique :

\begin{lstlisting}[language=C]
for (size_t i = 0; i < N; ++i) {
    P2[i] = calloc(M, sizeof(double));
}
\end{lstlisting}

La mémoire des lignes est sur le heap.
Le tableau des pointeurs est sur la stack.

\section{Cas 3 : double pointeur}

\begin{lstlisting}[language=C]
double** P3 = calloc(N, sizeof(double*));

for (size_t i = 0; i < N; ++i) {
    P3[i] = calloc(M, sizeof(double));
}
\end{lstlisting}

Ici :

\begin{itemize}
\item P3 est un pointeur vers pointeur.
\item La structure complète est sur le heap.
\item Les blocs ne sont pas nécessairement contigus.
\end{itemize}

\section{Différences cruciales}

\subsection{Continuité mémoire}

\begin{itemize}
\item P1 : mémoire entièrement continue.
\item P2 : N pointeurs contigus, lignes non contiguës.
\item P3 : tout dynamique, non garanti contigu.
\end{itemize}

\subsection{Allocation}

\begin{itemize}
\item P1 : allocation statique.
\item P2 : mix stack + heap.
\item P3 : entièrement heap.
\end{itemize}

\subsection{Dimensions variables}

\begin{itemize}
\item P1 : toutes les lignes ont M colonnes.
\item P2 et P3 : chaque ligne peut avoir une taille différente.
\end{itemize}

\section{Attention importante}

\[
\boxed{
T[N][M] \neq T** 
}
\]

Un tableau bidimensionnel n’est pas un double pointeur.

Un tableau multidimensionnel :

\begin{itemize}
\item est alloué de façon continue,
\item ressemble à un pointeur constant vers bloc continu.
\end{itemize}

Il ne faut jamais confondre :

\begin{lstlisting}[language=C]
int a[N][M];
\end{lstlisting}

avec :

\begin{lstlisting}[language=C]
int** a;
\end{lstlisting}

Ce sont des objets fondamentalement différents.

\section{Arithmétique des adresses}

Exemple :

\begin{lstlisting}[language=C]
printf("%ld\n", &P1[1][2] - &P1[0][0]);
\end{lstlisting}

Dans le cas de P1 :

\[
\text{distance} = 1 \times M + 2
\]

Car la mémoire est continue.

Dans le cas de P2 ou P3 :

\begin{itemize}
\item La distance peut être très grande.
\item Les blocs peuvent être éloignés en mémoire.
\item Stack et heap peuvent être très séparés.
\end{itemize}

\section{Résumé final}

\begin{itemize}
\item Un tableau 1D ressemble à un pointeur constant.
\item Un tableau 2D n’est PAS un double pointeur.
\item Les tableaux multidimensionnels sont continus.
\item Les doubles pointeurs ne garantissent aucune continuité.
\item Toujours distinguer allocation stack et heap.
\end{itemize}

\[
\boxed{
\text{Continuité mémoire = différence fondamentale}
}
\]

\section{Arithmétique des pointeurs}

L’arithmétique des pointeurs permet de se déplacer efficacement
dans une zone mémoire.

On peut utiliser :

\begin{itemize}
\item +, -
\item ++, --
\item différence de pointeurs
\end{itemize}

\section{Principe fondamental}

Si \texttt{p} est un pointeur sur un objet de type T :

\[
\boxed{
p + 1 \text{ pointe vers l’objet suivant de type T}
}
\]

Il ne faut pas penser en octets.

Il faut penser :

\[
\text{objet suivant}
\]

En réalité, le compilateur ajoute :

\[
sizeof(T)
\]

à l’adresse, mais il ne faut pas raisonner ainsi.

\section{Exemple : parcours d’un tableau}

\begin{lstlisting}[language=C]
int tab[N];

int* p = tab;
int* const end = tab + N;

while (p < end) {
    printf("%d\n", *p);
    ++p;
}
\end{lstlisting}

Analyse :

\begin{itemize}
\item tab est le début du tableau.
\item tab + N pointe juste après le dernier élément.
\item ++p avance d’un entier.
\end{itemize}

Important :

\[
\boxed{
* p \equiv p[0]
}
\]

\section{Pointeur constant}

\begin{lstlisting}[language=C]
const int* const end = tab + N;
\end{lstlisting}

Cela signifie :

\begin{itemize}
\item On ne modifie pas les valeurs pointées.
\item On ne modifie pas le pointeur lui-même.
\end{itemize}

\section{Exemple classique : parcours de chaîne}

Considérons :

\begin{lstlisting}[language=C]
char* s = "ABC";
char* p = s;
char c;
\end{lstlisting}

Boucle typique :

\begin{lstlisting}[language=C]
while ((c = *p++) != '\0') {
    putchar(c);
}
\end{lstlisting}

\section{Comprendre *p++}

Question :

Est-ce :

\begin{lstlisting}
(*p)++
\end{lstlisting}

ou

\begin{lstlisting}
*(p++)
\end{lstlisting}

Réponse :

\[
\boxed{
*p++ \equiv *(p++)
}
\]

Car l’opérateur postfixe ++ a une priorité plus élevée que *.

\section{Ordre des opérations}

\begin{lstlisting}[language=C]
c = *p++;
\end{lstlisting}

Équivalent à :

\begin{enumerate}
\item Lire la valeur pointée par p.
\item Affecter à c.
\item Incrémenter p.
\end{enumerate}

Donc :

\[
\text{lecture puis déplacement}
\]

\section{Différence avec (*p)++}

\begin{lstlisting}[language=C]
(*p)++;
\end{lstlisting}

Ici :

\begin{itemize}
\item On modifie la valeur pointée.
\item On n’avance pas le pointeur.
\end{itemize}

On incrémente dans l’espace des valeurs,
pas dans l’espace des pointeurs.

\section{Différence entre *p++ et *++p}

\begin{lstlisting}[language=C]
*p++
\end{lstlisting}

Lecture puis déplacement.

\begin{lstlisting}[language=C]
*++p
\end{lstlisting}

Déplacement puis lecture.

Ordre inverse.

\section{Pourquoi utiliser une variable temporaire ?}

Dans :

\begin{lstlisting}[language=C]
while ((c = *p++) != '\0') {
    ...
}
\end{lstlisting}

Si on écrivait simplement :

\begin{lstlisting}[language=C]
while (*p++ != '\0') {
    putchar(*p);   // ERREUR logique
}
\end{lstlisting}

On lirait un caractère en avance.

La variable temporaire garantit que
l’on travaille sur le bon caractère.

\section{Différence de pointeurs}

On peut soustraire deux pointeurs :

\begin{lstlisting}[language=C]
ptrdiff_t d = p2 - p1;
\end{lstlisting}

Le résultat :

\[
\boxed{
\text{nombre d’objets entre p1 et p2}
}
\]

Type retourné :

\[
\boxed{
ptrdiff_t
}
\]

Ce n’est pas un int.

Pourquoi ?

\begin{itemize}
\item Peut être négatif.
\item Doit pouvoir représenter de grandes distances mémoire.
\end{itemize}

\section{Indexation et équivalences}

\begin{lstlisting}[language=C]
tab[i]
\end{lstlisting}

est strictement équivalent à :

\begin{lstlisting}[language=C]
*(tab + i)
\end{lstlisting}

Et même :

\begin{lstlisting}[language=C]
i[tab]
\end{lstlisting}

Car :

\[
a[b] \equiv *(a + b)
\]

Même si cette écriture est déconseillée.

\section{Tableaux et passage aux fonctions}

Quand on écrit :

\begin{lstlisting}[language=C]
void f(int tab[]);
\end{lstlisting}

Le compilateur traduit en :

\begin{lstlisting}[language=C]
void f(int* tab);
\end{lstlisting}

La taille du tableau n’est pas connue.

\[
\boxed{
\text{Un tableau passé en argument devient un pointeur.}
}
\]

\section{Résumé}

\begin{itemize}
\item p + 1 signifie objet suivant.
\item *p++ = lecture puis déplacement.
\item (*p)++ = modification de la valeur.
\item p2 - p1 retourne un ptrdiff_t.
\item tab[i] = *(tab + i).
\item Les tableaux passés aux fonctions deviennent des pointeurs.
\end{itemize}

\[
\boxed{
\text{Toujours raisonner en objets, jamais en octets.}
}

\section{L’opérateur \texttt{sizeof}}

L’opérateur \texttt{sizeof} renvoie la taille en mémoire (en octets)
du type ou de l’expression donnée.

Syntaxe :

\begin{lstlisting}[language=C]
sizeof(type)
sizeof(expression)
\end{lstlisting}

\section{Point fondamental}

\[
\boxed{
\text{L’expression passée à sizeof n’est PAS évaluée.}
}
\]

Exemple piégeux :

\begin{lstlisting}[language=C]
int i = 10;
sizeof(i++);   // i n'est PAS incrementé
\end{lstlisting}

Ceci est strictement équivalent à :

\begin{lstlisting}[language=C]
sizeof(i);
\end{lstlisting}

\section{Exemples simples}

\begin{lstlisting}[language=C]
sizeof(double)   // taille d'un double
sizeof(int)      // taille d'un int
\end{lstlisting}

Si :

\begin{lstlisting}[language=C]
int tab[N];
\end{lstlisting}

Alors :

\begin{lstlisting}[language=C]
sizeof(tab)
\end{lstlisting}

renvoie :

\[
N \times sizeof(int)
\]

\section{Calcul du nombre d’éléments d’un tableau}

Idiom classique :

\begin{lstlisting}[language=C]
size_t n = sizeof(tab) / sizeof(tab[0]);
\end{lstlisting}

Cela fonctionne car :

\[
\frac{N \times sizeof(int)}{sizeof(int)} = N
\]

\section{Attention : tableau vs pointeur}

Considérons :

\begin{lstlisting}[language=C]
int tab[1000];
int* t = tab;
\end{lstlisting}

Alors :

\begin{lstlisting}[language=C]
sizeof(tab)   // 1000 * sizeof(int)
sizeof(t)     // sizeof(int*)
\end{lstlisting}

\[
\boxed{
sizeof(t) \neq sizeof(tab)
}
\]

\section{Erreur classique}

Dans une fonction :

\begin{lstlisting}[language=C]
void f(int t[]) {
    size_t n = sizeof(t) / sizeof(t[0]);
}
\end{lstlisting}

Ceci est faux.

Pourquoi ?

Parce que :

\begin{lstlisting}
int t[]
\end{lstlisting}

est transformé par le compilateur en :

\begin{lstlisting}
int* t
\end{lstlisting}

Donc :

\[
sizeof(t) = sizeof(int*)
\]

Le résultat sera 1, 2 ou autre petite valeur,
mais certainement pas la taille réelle du tableau.

\section{Règle absolue}

\[
\boxed{
Un tableau ne connaît jamais sa taille.
}
\]

Toujours passer la taille en argument :

\begin{lstlisting}[language=C]
void f(int t[], size_t n);
\end{lstlisting}

\section{Bug classique avec l’arithmétique}

Erreur fréquente :

\begin{lstlisting}[language=C]
int tab[1000];
int* end = tab + sizeof(tab);
\end{lstlisting}

Problème :

\[
sizeof(tab) = 1000 \times sizeof(int)
\]

Donc on avance beaucoup trop loin.

Correct :

\begin{lstlisting}[language=C]
int* end = tab + sizeof(tab)/sizeof(tab[0]);
\end{lstlisting}

\section{Autre bug typique}

\begin{lstlisting}[language=C]
for (int* p = tab; p < tab + sizeof(tab); ++p) {
    ...
}
\end{lstlisting}

Ici, on dépasse largement la mémoire du tableau.

Cela peut provoquer :

\begin{itemize}
\item lecture hors limites,
\item écriture hors limites,
\item buffer overflow.
\end{itemize}

\section{Bonne pratique}

Toujours écrire :

\begin{lstlisting}[language=C]
size_t n = sizeof(tab)/sizeof(tab[0]);

for (int* p = tab; p < tab + n; ++p) {
    ...
}
\end{lstlisting}

\section{Résumé}

\begin{itemize}
\item sizeof renvoie une taille en octets.
\item L’expression n’est jamais évaluée.
\item sizeof(tab) fonctionne uniquement si tab est un vrai tableau.
\item Dans une fonction, un tableau devient un pointeur.
\item Toujours passer explicitement la taille.
\item Attention aux dépassements lors de tab + sizeof(tab).
\end{itemize}

\[
\boxed{
\text{sizeof est simple, mais source de nombreux bugs subtils.}
}

\section{Bonus : Flexible Array Member}

Le \textbf{Flexible Array Member} (FAM) est un pattern avancé,
souvent utilisé en programmation système.

L’idée est de créer une structure contenant :

\begin{itemize}
\item une taille,
\item un tableau dynamique stocké directement à la suite en mémoire.
\end{itemize}

\section{Motivation}

Classiquement, on écrirait :

\begin{lstlisting}[language=C]
typedef struct {
    size_t size;
    double* data;
} vector;
\end{lstlisting}

Dans ce cas :

\begin{itemize}
\item La structure contient un pointeur.
\item Les données sont ailleurs dans le heap.
\item Il y a une indirection supplémentaire.
\end{itemize}

Le Flexible Array Member permet d’éviter cela.

\section{Définition}

\begin{lstlisting}[language=C]
typedef struct {
    size_t size;
    double data[];
} vector;
\end{lstlisting}

Important :

\[
\boxed{
Le tableau flexible doit être le dernier champ de la structure.
}
\]

Il n’a pas de taille déclarée.

\section{Principe mémoire}

La structure contient :

\begin{itemize}
\item la partie fixe (size),
\item puis une zone contiguë pour les données.
\end{itemize}

On alloue alors :

\[
sizeof(vector) + (n - 1) \times sizeof(double)
\]

Pourquoi \(n - 1\) ?

Parce que \texttt{sizeof(vector)} inclut déjà la place
pour un élément de \texttt{data} (selon l’ancienne convention
avec data[1]).  
Avec la forme moderne data[], il faut simplement
allouer l’espace total nécessaire.

\section{Allocation sécurisée}

On doit éviter les overflows lors du calcul.

\begin{lstlisting}[language=C]
vector* create_vector(size_t n) {
    if (n == 0) return NULL;

    if (n > (SIZE_MAX - sizeof(vector)) / sizeof(double)) {
        return NULL;
    }

    vector* v = malloc(sizeof(vector) +
                       n * sizeof(double));

    if (v == NULL) return NULL;

    v->size = n;

    return v;
}
\end{lstlisting}

\section{Utilisation}

Une fois alloué :

\begin{lstlisting}[language=C]
vector* v = create_vector(10);

v->data[0] = 1.0;
v->data[1] = 2.0;
\end{lstlisting}

Les données sont stockées immédiatement après la structure
en mémoire, de manière contiguë.

\section{Avantage principal}

Comparons les deux approches.

\subsection{Avec pointeur}

\begin{lstlisting}[language=C]
typedef struct {
    size_t size;
    double* data;
} vector;
\end{lstlisting}

Mémoire :

\begin{itemize}
\item structure sur heap,
\item tableau séparé sur heap,
\item une indirection supplémentaire.
\end{itemize}

\subsection{Avec flexible array member}

\begin{lstlisting}[language=C]
typedef struct {
    size_t size;
    double data[];
} vector;
\end{lstlisting}

Mémoire :

\begin{itemize}
\item une seule allocation,
\item données contiguës,
\item pas d’indirection,
\item meilleure localité cache.
\end{itemize}

\section{Schéma conceptuel}

Structure classique :

\[
[ size ][ ptr ] \rightarrow [ data... ]
\]

Flexible array :

\[
[ size ][ data... data... data... ]
\]

\section{Limites}

\begin{itemize}
\item Pattern plus avancé.
\item Doit être alloué dynamiquement.
\item Le champ flexible doit être en dernier.
\item Impossible d’avoir d’autres champs après.
\end{itemize}

\section{Pourquoi c’est utilisé}

\begin{itemize}
\item Code système.
\item Structures variables.
\item Optimisation mémoire.
\item Protocoles réseau.
\end{itemize}

\section{Résumé}

\begin{itemize}
\item Permet d’avoir une structure + données contiguës.
\item Évite une indirection.
\item Améliore la localité mémoire.
\item Nécessite une allocation calculée avec soin.
\end{itemize}

\[
\boxed{
\text{Flexible Array Member = structure extensible en mémoire contiguë.}
}
